\input{preamble}

\begin{document}

\title{K3 surfaces exam}
\author{Daniel González Casanova Azuela}
\maketitle

\section*{Complex curves on a K3 surface}
\label{section-curves}

\begin{exercise}[2.6a]
\label{exercise-deformations-curve-in-k3}
Let $C$ be a smooth genus $g$ curve which
can be embedded in a K3 surface $M$, and
$X$
the family of all deformations of $C$ in
$M$.
Prove that $\dim X \leq g$.
\end{exercise}

\begin{proof}
We are asked to compute
$\dim X = \dim|C|=h^0(\mathcal{O}_S(C))-1$.

By definition of Euler characteristic,
$$
\chi(\mathcal{O}(C))
=
h^0(\mathcal{O}(C))
- h^1(\mathcal{O}(C))
+ h^2(\mathcal{O}(C)).
$$
To compute this Euler characteristic we use
Riemann-Roch theorem for curves in surfaces:
$$
\chi(\mathcal{O}(C))
=
\chi(\mathcal{O}_S)
+\frac{1}{2}(\mathcal{O}(C),\mathcal{O}(C)-K_S)
=
2+\frac{1}{2}\mathcal{O}(C)^2
$$
since $S$ is a K3 surface.

To compute the self-intersection number
we use genus formula:
$$
(\mathcal{O}(C).\mathcal{O}(C)+K_S)
=
\mathcal{O}(C)^2
=
2g-2
$$
since $S$ is a K3 surface.

Thus $\chi(\mathcal{O}(C))=g+1$ and
$$
g+1=h^0(\mathcal{O}(C))
-h^1(\mathcal{O}(C))
+h^2(\mathcal{O}(C)).
$$

By Serre duality and $S$ being a K3,
we obtain two equations:
$$
h^2(\mathcal{O}(C))
=
h^0(\mathcal{O}(-C)),
\qquad 
h^1(\mathcal{O}(C))
=
h^1(\mathcal{O}(-C)).
$$

The ideal sheaf $\mathcal{O}(-C)$
exact sequence gives
the long exact sequence
$$
0 \to H^0(\mathcal{O}(-C))
\hookrightarrow
H^0(\mathcal{O}_S)=\mathbb{C}
\to
H^0(\mathcal{O}_C)=\mathbb{C}
\to 
H^1(\mathcal{O}(-C))
\to
\cdots
$$
First note that
$h^0(-C) = 0$
since the global sections of the
ideal sheaf $\mathcal{O}(-C)$
are global holomorphic functions on $S$
vanishing along $C$---only the
constant zero function lies here.
Thus
$H^0(\mathcal{O}_S) =\mathbb{C}
\to
H^0(\mathcal{O}_C)=\mathbb{C}$
in the long exact sequence is an injection,
and thus surjective, making
$H^1(\mathcal{O}(-C))=0$.
Thus $\dim X = g$.
\end{proof}

\begin{exercise}[2.7]
\label{exercise-singular-curve-on-k3}
Let $C$ be a singular, irreducible complex
curve on a K3 surface $X$. Prove that
$(C.C)\geq 0$.
\end{exercise}

\begin{proof}
This is an application of the genus
formula. Since
$K_X=\mathcal{O}_X$ it suffices to show that
the arithmetic genus
$p_a(C)$ is not zero. This follows from the
fact that $C$ is singular, since its
arithmetic genus, defined as 
$p_a(C):=1-\chi(\mathcal{O}_C)$,
satisfies
\begin{equation}
\label{equation-normalization-genus}
p_a(C)=g(C) + h^0(\delta)
\end{equation}
where $g(C)$ is the geometric genus of $C$,
defined as
the genus of the normalization $\tilde{C}$,
i.e. $g(C):=g(\tilde{C})$,
while $\delta:=\nu_*\mathcal{O}_{\tilde{C}}
/\mathcal{O}_C$
where $\nu:\tilde{C} \to C$ is the
normalization map.

To prove Equation
\ref{equation-normalization-genus}
first note that
$p_a(C)=h^1(\mathcal{O}_C)$
and $g(C) =
h^1(\mathcal{O}_{\widetilde{C}})$,
where the latter uses Serre duality since
normalization of a curve is smooth.

Now use the normalization
long exact sequence in cohomology
$$
0 \to
H^0(\mathcal{O}_C)\to
H^0(\nu_*\mathcal{O}_{\widetilde{C}})\to
H^0(\delta)\to
H^1(\mathcal{O}_C)\to
H^1(\nu_*\mathcal{O}_{\widetilde{C}})\to
H^1(\delta)
$$
Note the following three facts:
the map
$H^0(\mathcal{O}_C)
\to
H^0(\nu_*\mathcal{O}_{\widetilde{C}})$
is an isomorphism,
$H^1(\nu_*\mathcal{O}_{\widetilde{C}})
\cong H^1(\mathcal{O}_{\widetilde{C}})$
and
$H^1(\delta)=0$.
The first is due to both sheaves'
global sections being
constant functions,
the second by $\nu$ being finite
(Lecture 9),
and the third because $\delta$
is supported in finitely many points.
The latter is because the support of
$\delta$ are the singularities of $C$ 
(which must be finitely many
by compactness): the stalk
$\delta_p$ is the quotient of the
integral closure of the stalk
$\mathcal{O}_p$ by $\mathcal{O}_p$,
which is nonzero if and only if $p$
is not a smooth point.
Thus we have the short exact sequence
$$
\xymatrix{
0\ar[r]
&H^0(\delta)\ar[r]
&H^1(\mathcal{O}_C)\ar[r]
&H^1(\nu_*\mathcal{O}_{\widetilde{C}})\ar[r]
&0
}
$$
yielding
Equation \ref{equation-normalization-genus}.

Finally, since $C$ is singular
$\delta$ is nonzero and thus
$h^0(\delta)>0$.
\end{proof}

\begin{lemma}[Genus formula]
\label{lemma-genus-formula}
Let $C$ be any curve on a surface $S$. Then
\begin{equation}
\label{equation-genus-formula}
2p_a(C)-2=(C.C+K_S).
\end{equation}
\end{lemma}

\begin{proof}
By Remark
\ref{remark-euler-characteristic-additive}
below we can take Euler characteristic
of the ideal sheaf exact sequence of $C$
to obtain
$$
\chi(\mathcal{O}(-C))+\chi(\mathcal{O}_C)=
\chi(\mathcal{O}_S).
$$
Since $p_a(C)=h^0(K_S)
=1-\chi(\mathcal{O}_C)$, we
obtain that
$p_a(C)=1-\chi(\mathcal{O}_S)-
\chi(\mathcal{O}(-C))$.
To compute $\chi(\mathcal{O}(-C))$ we use
Riemann-Roch formula for curves on surfaces
to get
$$
\chi(\mathcal{O}(-C))=\chi(\mathcal{O}_S)
+\frac{(\mathcal{O}(-C).
\mathcal{O}(-C)-K_S)}{2}
$$
Taking duals on both entries of the
intersection form and substituting in our
previous expression for $p_a(C)$ we obtain
the result.
\end{proof}

\begin{remark}
\label{remark-euler-characteristic-additive}
To prove that Euler characteristic is
additive, i.e. that for an exact sequence
$$
\xymatrix{
0\ar[r]
&E\ar[r]
&F\ar[r]
&G\ar[r]
&0
}
$$
we must have
$$
\chi(F)=\chi(E)+\chi(G),
$$
we take the long exact sequence in cohomology
and use the general fact that for
a finite exact sequence of vector spaces
$$
\xymatrix{
0\ar[r]
&V_0\ar[r]
&V_1\ar[r]
&\cdots\ar[r]
&V_n\ar[r]
&0
}
$$
we must have
$$
\sum_i(-1)^i\dim V_i,
$$
which basically follows by dimension theorem.
\end{remark}

\end{document}
